\chapter{Testing and Evaluation}
\label{ch:testing-evaluation}

We tested the design by running each configuration on a set of
hand-written parsers and reviewing the synthesised state machine for
correctness. A small Python helper renders the synthesised state
machine as a PDF, which made the manual review easier.

We evaluate the design separately for input programs that use only
exact comparisons and for programs that involve ternary
comparisons since the supported configurations differ in the two
cases. 

\section{Exact comparisons}
\label{sec:eval-exact}

We ran four benchmarks with each of the following configurations
for both Design~1 and Design~2. \texttt{max\_states} and
\texttt{max\_table\_entries} were both set to ten, and a timeout of
20 minutes was imposed on each run.

\begin{enumerate}
    \item Constant synthesis and field minimisation.
    \item Field minimisation only.
    \item Constant synthesis only.
    \item No optimisation.
\end{enumerate}

Table~\ref{tab:eval-exact-fieldsizes} reports the maximum bit width
of a packet with and without field minimisation on these benchmarks.
Table~\ref{tab:exact-results} reports the resulting state count,
table-entry count and synthesis time for each configuration for
both designs.

\begin{table}[H]
\centering
\rowcolors{2}{blue!5}{white}
\begin{tabular}{lcc}
\hline
\rowcolor{blue!20}
\textbf{Benchmark} & \textbf{field min.\ on (bits)} & \textbf{field min.\ off (bits)} \\
\hline
\texttt{linear}                & 3 & 72  \\
\texttt{single\_if}            & 4 & 56  \\
\texttt{multiple\_if}          & 7 & 88  \\
\texttt{check\_after\_extract} & 5 & 72  \\
\hline
\end{tabular}
\caption{Maximum bit width of a packet with and without
field minimisation, on the exact-comparison benchmarks.}
\label{tab:eval-exact-fieldsizes}
\end{table}

\begin{table}[H]
\centering
\resizebox{\textwidth}{!}{%
\begin{tabular}{l|cc|ccc|ccc}
\hline
\rowcolor{blue!20}
\textbf{Benchmark} & \textbf{CS} & \textbf{FM}
& \multicolumn{3}{c|}{\textbf{Design 1}}
& \multicolumn{3}{c}{\textbf{Design 2}} \\
\rowcolor{blue!20}
 & & & \textbf{States} & \textbf{Entries} & \textbf{Time (s)}
     & \textbf{States} & \textbf{Entries} & \textbf{Time (s)} \\
\hline
\texttt{linear}                 & on  & off & 3 & 0 & 1.23   & 3 & 0 & 2.44    \\
\texttt{linear}                 & on  & on  & 3 & 0 & 0.43   & 3 & 0 & 0.24    \\
\texttt{linear}                 & off & off & 3 & 0 & 1.44   & 3 & 0 & 1.84    \\
\texttt{linear}                 & off & on  & 3 & 0 & 0.45   & 3 & 0 & 0.19    \\
\hline
\texttt{single\_if}             & on  & off & 3 & 1 & 68.08  & 3 & 1 & 51.89   \\
\texttt{single\_if}             & on  & on  & 3 & 1 & 14.86  & 3 & 1 & 21.98   \\
\texttt{single\_if}             & off & off & 3 & 1 & 55.11  & 3 & 1 & 50.49   \\
\texttt{single\_if}             & off & on  & 3 & 1 & 11.18  & 3 & 1 & 23.41   \\
\hline
\texttt{check\_after\_extract} & on  & off & 4 & 2 & 132.97 & 4 & 2 & 164.23  \\
\texttt{check\_after\_extract} & on  & on  & 4 & 2 & 35.47  & 4 & 2 & 144.92  \\
\texttt{check\_after\_extract} & off & off & 4 & 2 & 159.22 & 4 & 2 & 279.61  \\
\texttt{check\_after\_extract} & off & on  & 4 & 2 & 39.28  & 4 & 2 & 133.84  \\
\hline
\texttt{multiple\_if}           & on  & off & 7 & 3 & 931.27 & 6 & 2 & 557.88  \\
\texttt{multiple\_if}           & on  & on  & 7 & 3 & 339.76 & 6 & 2 & 249.91  \\
\texttt{multiple\_if}           & off & off & 7 & 3 & 696.18 & 6 & 2 & 652.19  \\
\texttt{multiple\_if}           & off & on  & 7 & 3 & 421.73 & 6 & 2 & 330.57  \\
\hline
\end{tabular}%
}
\caption{Synthesis results on the exact-comparison benchmarks
across all four \textsf{CS}/\textsf{FM} configurations, for both
designs. \textsf{CS} stands for constant synthesis and \textsf{FM}
for field minimisation.}
\label{tab:exact-results}
\end{table}

Field minimisation was the most effective optimisation. It produced
a 3--5$\times$ speedup over the unminimised baseline on most
benchmarks. Constant synthesis produced a reasonable speedup on
some benchmarks but performed worse on others. On
\texttt{multiple\_if}, the second design produced an FSM with one
fewer state and one fewer table entry than the first
and was faster. On the other benchmarks, where both
designs produced state machines of the same size, the first design
ran faster due to its smaller constraints.

We additionally evaluated the best-performing configuration on
realistic parsers translated from the ParserHawk benchmarks, with
a timeout of 30 minutes, a maximum of 10 states and a maximum of
10 table entries. Synthesis did not complete within the allotted
time. This indicates that more efficient optimisation techniques
are required to make the approach practically applicable, and is
an important direction for future work.

\clearpage
\section{Ternary comparisons}
\label{sec:eval-ternary}

Only Design~2 supports ternary comparisons, as discussed in
Section~\ref{sec:synthesis}, so we evaluated it with two
configurations: field minimisation on, and field minimisation off.
Field minimisation was applied only on the fields that participate
in exact comparisons or in no comparison at all. Fields involved in
ternary comparisons were left untouched.

A timeout of 20 minutes was imposed on each run. The benchmarks
were run with eight states and eight table entries as the maximum.
Table~\ref{tab:eval-ternary-fieldsizes} reports the maximum bit width of a packet.
Table~\ref{tab:eval-ternary} reports the resulting state count,
table-entry count and synthesis time. A dash indicates that the
configuration timed out.

\begin{table}[H]
\centering
\rowcolors{2}{blue!5}{white}
\begin{tabular}{lcc}
\hline
\rowcolor{blue!20}
\textbf{Benchmark} & \textbf{field min.\ on (bits)} & \textbf{field min.\ off (bits)} \\
\hline
\texttt{single\_if\_else}      & 6  & 12  \\
\texttt{multiple\_if}          & 9  & 36  \\
\texttt{check\_after\_extract} & 10 & 72  \\
\texttt{nested\_if\_else2}     & 6  & 12  \\
\texttt{nested\_if\_else1}     & 12 & 104 \\
\hline
\end{tabular}
\caption{Maximum bit width of a packet with and without
field minimisation, on the ternary benchmarks.}
\label{tab:eval-ternary-fieldsizes}
\end{table}

\begin{table}[H]
\centering
\rowcolors{2}{blue!5}{white}
\begin{tabular}{lcccc}
\hline
\rowcolor{blue!20}
\textbf{Benchmark} & \textbf{Field min.} & \textbf{States} & \textbf{Entries} & \textbf{Time (s)} \\
\hline
\texttt{single\_if\_else}      & on  & 4 & 1 & 18.85   \\
\texttt{single\_if\_else}      & off & 4 & 1 & 18.55   \\
\texttt{multiple\_if}          & on  & 4 & 3 & 432.67  \\
\texttt{multiple\_if}          & off & 4 & 3 & 484.76  \\
\texttt{check\_after\_extract} & on  & 4 & 2 & 79.64   \\
\texttt{check\_after\_extract} & off & 4 & 2 & 206.32  \\
\texttt{nested\_if\_else2}     & on  & 4 & 2 & 56.29   \\
\texttt{nested\_if\_else2}     & off & 4 & 2 & 64.38   \\
\texttt{nested\_if\_else1}     & on  & 6 & 3 & 602.07  \\
\texttt{nested\_if\_else1}     & off & - & - & timeout \\
\hline
\end{tabular}
\caption{Synthesis results on ternary benchmarks with Design~2.
Field minimisation was applied only on fields that do not
participate in ternary comparisons.}
\label{tab:eval-ternary}
\end{table}

The state and table-entry counts match across the two
configurations for every benchmark that completed and the run with 
field minimisation on runs faster always.
The speedup from field minimisation is substantial
on \texttt{check\_after\_extract} and \texttt{nested\_if\_else1}.
On \texttt{nested\_if\_else1} the minimised configuration completes
in about ten minutes while the unminimised one times out.

